\documentclass[runningheads]{llncs}

\usepackage[T1]{fontenc}
\usepackage{graphicx}
\usepackage{amsmath}
\usepackage{amssymb}
\usepackage{booktabs}
\usepackage{array}
\usepackage{url}
\usepackage{color}
\usepackage{hyperref}
\usepackage{microtype}
\renewcommand\UrlFont{\color{blue}\rmfamily}

\title{Explainable Quantum-Inspired Multi-label Classification of Non-Functional Requirements}
\titlerunning{Explainable Quantum-Inspired Multi-label NFR Classification}

\author{Nguyen Hoang Phuc\inst{1} \and Le Doan Gia Hung\inst{1} \and Lai Duc Hung\inst{1}}
\authorrunning{N. H. Phuc et al.}

\institute{FPT University, Ho Chi Minh City, Vietnam \\
\email{nphuc4748@gmail.com, lhung291005@gmail.com, hungld@fpt.edu.vn}}

\begin{document}
\maketitle

\begin{abstract}
Non-functional requirements (NFRs) frequently overlap across quality attributes
such as security, performance, usability, and availability, which makes NFR
classification a natural multi-label problem. This paper presents an explainable
quantum-inspired approach for multi-label NFR classification. Requirements are
represented as normalized semantic states, NFR labels as projection directions,
and label scores as squared projection amplitudes; a label interference matrix
models NFR co-occurrence, and token-level amplitude contributions provide
interpretability. We evaluate several quantum-inspired variants on the NICE
multi-label dataset against TF-IDF Logistic Regression, TF-IDF Linear SVM, a
Sentence-BERT baseline, and a label-frequency baseline. Under 5-fold
cross-validation, Sentence-BERT achieves the highest Macro-F1, while classical
TF-IDF baselines and the hybrid quantum--SVM model perform comparably. The
Macro-F1 differences between the hybrid model and TF-IDF baselines are not
statistically significant in a Wilcoxon signed-rank test. The results indicate
that the quantum-inspired projection is competitive with, rather than superior
to, strong classical baselines, while offering an interpretable, geometry-based
account of overlapping NFR semantics. We further provide a deletion-based
evaluation of the token-level explanations.
\keywords{Requirements Engineering \and Non-Functional Requirements \and Multi-label Classification \and Quantum-inspired Model \and Explainable AI \and Software Engineering}
\end{abstract}

\section{Introduction}
Requirements Engineering (RE) is a critical activity in software development
because requirement quality strongly affects downstream architecture,
implementation, testing, and maintenance decisions. Among requirement types,
non-functional requirements (NFRs) are particularly difficult to identify and
classify automatically. NFRs describe quality attributes such as security,
performance, usability, availability, scalability, and maintainability. Unlike
many functional requirements, NFRs are often ambiguous, context-dependent, and
overlapping.

For example, a requirement stating that payment information must be encrypted
and accessible only to authorized administrators clearly relates to security. If
the same requirement also constrains availability, auditability, or response
time, it may express multiple NFR categories at once. Therefore, NFR
classification is better treated as a multi-label problem rather than a strict
single-label problem \cite{tsoumakas2007multilabel,zhang2014review}.

This paper investigates the research topic: \emph{Explainable Quantum-Inspired
Multi-label Classification of Non-Functional Requirements}. The key idea is to
represent each requirement as a semantic state and estimate its degree of
association with multiple NFR labels using projection amplitudes. The approach
is quantum-inspired rather than quantum-hardware-based: it runs on ordinary
computers and borrows mathematical intuitions such as state representation,
projection amplitude, and interference from geometric information retrieval and
quantum language modeling \cite{vanrijsbergen2004geometry,sordoni2013quantum,li2019cnm}.

The basic projection component is closely related to centroid-based text
classification. The contribution of this work is therefore not a claim of raw
performance superiority, but a combination of (i) an interpretable geometric
scoring view, (ii) an interference adjustment for label co-occurrence, and
(iii) hybrid calibration with a discriminative classifier.

The contributions of this paper are:
\begin{itemize}
    \item a reproducible implementation for multi-label NFR classification experiments;
    \item quantum-inspired projection, contrastive projection, and hybrid quantum--SVM variants;
    \item token-level semantic amplitude explanations and a deletion-based faithfulness test;
    \item empirical comparison against TF-IDF Logistic Regression, TF-IDF Linear SVM, Sentence-BERT, and a label-frequency baseline on the NICE dataset;
    \item an ablation study of projection type, label interference, and hybrid fusion weight.
\end{itemize}

\section{Related Work}
Automated NFR classification has been studied for many years in Requirements
Engineering. Early work by Cleland-Huang et al. explored automated
classification of NFRs such as security, performance, and usability
\cite{clelandhuang2007automated}. Later studies introduced supervised machine
learning and dependency-based explanations for requirements classification
\cite{kurtanovic2017automatically,dalpiaz2019requirements}. The PROMISE and
PROMISE-expanded datasets have supported repeatable experiments on requirement
classification \cite{mitrevski_promise_exp}.

NICE introduced a multi-label NFR dataset based on the augmented PROMISE dataset
and relabeled NFR sub-classes such as availability, security, and performance
using a multi-label schema \cite{nice_dataset,nice_preprint}. This dataset is
especially relevant to the present work because the research problem is
explicitly multi-label. A direct numeric comparison with the NICE paper is
limited by differences in label filtering, splits, and evaluation protocol, so
we report our own baselines on the same downloaded data.

Traditional text classification commonly relies on bag-of-words or TF-IDF
features combined with classifiers such as Logistic Regression or Support Vector
Machines \cite{sebastiani2002machine}. Transformer-based models have become
strong baselines for requirements classification, including BERT-style transfer
learning approaches such as NoRBERT \cite{devlin2019bert,hey2020norbert}.
Sentence-BERT provides efficient sentence embeddings for downstream classifiers
\cite{reimers2019sentencebert}, and is therefore included as a stronger
embedding baseline in this study.

The quantum-inspired component is related to geometric information retrieval and
quantum-inspired language models, where term dependencies and matching can be
modeled through vector spaces, projections, and complex-valued or
wave-function-inspired representations
\cite{vanrijsbergen2004geometry,sordoni2013quantum,zhang2018quantum,li2019cnm}.
For explainability, this work is positioned near model explanation methods such
as LIME and SHAP \cite{ribeiro2016lime,lundberg2017shap}, and evaluates
faithfulness using a deletion test inspired by rationale evaluation work such
as ERASER \cite{deyoung2020eraser}.

\section{Research Questions}
This study is guided by three research questions:

\textbf{RQ1.} Can a quantum-inspired semantic projection model classify
multi-label NFRs competitively compared with classical and embedding baselines?

\textbf{RQ2.} Does the model provide useful explainability through token-level
semantic amplitude contributions?

\textbf{RQ3.} Is the observed behavior stable under cross-validation rather than
only in a single train/test split?

\section{Proposed Method}
\subsection{Requirement Representation}
Each requirement text is first converted into a TF-IDF vector:
\begin{equation}
    \mathbf{x}_i \in \mathbb{R}^{d},
\end{equation}
where $d$ is the vocabulary size. The vector is normalized to form a semantic
state:
\begin{equation}
    \lvert r_i \rangle = \frac{\mathbf{x}_i}{\left\|\mathbf{x}_i\right\|_2}.
\end{equation}
This normalization allows the requirement to be interpreted as a direction in
semantic space.

\subsection{Label Projection}
For each NFR label $c$, a label basis vector is estimated from the centroid of
positive training examples:
\begin{equation}
    \lvert c \rangle = \frac{1}{N_c}\sum_{i:y_{ic}=1}\lvert r_i \rangle,
\end{equation}
followed by normalization. The raw association between requirement $r_i$ and
label $c$ is computed as a squared projection amplitude:
\begin{equation}
    s_c(r_i) = \left(\langle c \mid r_i \rangle\right)^2.
\end{equation}
We call $s_c(r_i)$ an association score rather than a probability because the
scores over labels are not constrained to sum to one.

\subsection{Contrastive Label Projection}
The positive-only centroid projection is simple and interpretable, but it does
not explicitly model how a label differs from other labels. Therefore, this
work also evaluates a contrastive quantum-inspired variant. For each label $c$,
the contrastive direction is computed as:
\begin{equation}
    \lvert c_{\Delta} \rangle =
    \frac{\mu_c^{+} - \mu_c^{-}}
    {\left\|\mu_c^{+} - \mu_c^{-}\right\|_2},
\end{equation}
where $\mu_c^{+}$ is the centroid of requirements that contain label $c$, and
$\mu_c^{-}$ is the centroid of requirements that do not contain label $c$.

\subsection{Interference Adjustment}
To reflect label co-occurrence, the model estimates a label interference matrix
from training labels. Let $N_{kl}$ be the number of training samples containing
both labels $k$ and $l$, and let $N_k$ be the number of samples containing label
$k$. The normalized co-occurrence matrix is:
\begin{equation}
    M_{kl} = \frac{N_{kl}}{N_k},
\end{equation}
with the diagonal set to 1. The implementation mixes this matrix with the
identity matrix:
\begin{equation}
    \mathbf{M}' = (1-\lambda)\mathbf{I} + \lambda\mathbf{M},
\end{equation}
where $\lambda$ controls interference strength. For projection scores
$\mathbf{p}_i$, adjusted scores are computed as:
\begin{equation}
    \mathbf{s}_i = \mathrm{normalize}(\mathbf{p}_i \mathbf{M}').
\end{equation}

\subsection{Hybrid Quantum--SVM Fusion}
To improve stability on small and imbalanced datasets, a hybrid model combines
the contrastive quantum-inspired score with a one-vs-rest Linear SVM score:
\begin{equation}
    \mathbf{h}_i = \alpha \mathbf{q}_i + (1-\alpha)\mathbf{v}_i,
\end{equation}
where $\mathbf{q}_i$ is the contrastive quantum-inspired score vector,
$\mathbf{v}_i$ is the calibrated SVM score vector, and $\alpha$ is the quantum
fusion weight. In the final experiments, $\alpha=0.15$ was selected based on
validation behavior.

\subsection{Explainability}
For a predicted label $c$, token-level contribution is estimated from the
element-wise product between the requirement state and the label basis:
\begin{equation}
    e_{ijc} = r_{ij}c_j.
\end{equation}
Terms with the largest absolute contributions are reported as explanations. We
evaluate these explanations using a deletion test: if the highlighted terms are
faithful, deleting them should reduce the model score more than deleting random
nonzero terms.

\section{Experimental Setup}
\subsection{Datasets}
\textbf{NICE multi-label dataset.} The main experiment uses the NICE dataset
from Zenodo \cite{nice_dataset}. The downloaded CSV contains 622 requirements
and binary label columns for NFR classes. After filtering rows with at least one
NFR sub-class label and removing the empty ``Other'' label, the experiment uses
381 requirements and 11 NFR labels. Among these, 125 requirements contain more
than one NFR label, giving a label cardinality of 1.3648.

\textbf{PROMISE-expanded dataset.} A secondary experiment uses
PROMISE-expanded from the public requirements classification repository
\cite{mitrevski_promise_exp}. This dataset is treated as an 11-class
single-label NFR subtype classification task and is used as an auxiliary sanity
check rather than the main multi-label experiment.

\subsection{Compared Models}
The compared baselines are:
\begin{itemize}
    \item \textbf{Label-frequency baseline:} predicts labels according to training label frequencies.
    \item \textbf{TF-IDF Logistic Regression:} one-vs-rest Logistic Regression using unigram and bigram TF-IDF features.
    \item \textbf{TF-IDF Linear SVM:} one-vs-rest Linear SVM using unigram and bigram TF-IDF features.
    \item \textbf{Sentence-BERT + LR:} all-MiniLM-L6-v2 embeddings followed by one-vs-rest Logistic Regression.
\end{itemize}

The quantum-inspired models are:
\begin{itemize}
    \item \textbf{Quantum-inspired Projection:} positive-centroid semantic projection with label interference.
    \item \textbf{Contrastive Quantum Projection:} positive-minus-negative semantic projection with label interference.
    \item \textbf{Hybrid Quantum--SVM Fusion:} score-level fusion of contrastive quantum projection and Linear SVM, using a quantum fusion weight of 0.15.
\end{itemize}

\subsection{Metrics and Threshold Selection}
The main reported metrics are Micro-F1, Macro-F1, Hamming loss, and Label
Ranking Average Precision (LRAP). Macro-F1 is emphasized because the NFR labels
are imbalanced. Two threshold protocols are evaluated. First, a single global
threshold is selected on an internal validation split by maximizing Macro-F1
over thresholds from 0.05 to 0.95. Second, a per-label threshold protocol
selects an independent threshold for each NFR label on the validation split.

\section{Results}
\subsection{Single Train/Validation/Test Split on NICE}
Table \ref{tab:nice-single} reports results from a single
train/validation/test split on the NICE dataset.

\begin{table}[htbp]
\centering
\caption{NICE multi-label results on a single train/validation/test split, comparing label-frequency, TF-IDF, Sentence-BERT, and quantum-inspired models across Micro-F1, Macro-F1, Hamming loss, and LRAP. Sentence-BERT has the highest Macro-F1, while the hybrid quantum--SVM model is close in Macro-F1 but lower in LRAP.}
\label{tab:nice-single}
\resizebox{\textwidth}{!}{
\begin{tabular}{lccccc}
\toprule
Model & Threshold & Micro-F1 & Macro-F1 & Hamming Loss & LRAP \\
\midrule
Label-frequency baseline & 0.05 & 0.2221 & 0.2166 & 0.8751 & 0.4183 \\
TF-IDF Logistic Regression & 0.45 & 0.6299 & 0.5266 & 0.0901 & 0.7861 \\
TF-IDF Linear SVM & 0.45 & 0.6212 & 0.5205 & \textbf{0.0877} & 0.7843 \\
Sentence-BERT + LR & 0.50 & \textbf{0.6667} & \textbf{0.6049} & 0.0901 & \textbf{0.8286} \\
Quantum-inspired Projection & 0.85 & 0.6076 & 0.5618 & 0.0980 & 0.7694 \\
Contrastive Quantum Projection & 0.60 & 0.5848 & 0.5367 & 0.1123 & 0.7600 \\
Hybrid Quantum--SVM Fusion & 0.45 & 0.6347 & 0.6019 & 0.0964 & 0.7847 \\
\bottomrule
\end{tabular}
}
\end{table}

On this single split, Sentence-BERT obtains the highest Macro-F1 and LRAP. The
hybrid quantum--SVM model is close in Macro-F1, but this single split contains
only 381 usable instances and should not be read as evidence of superiority.

\subsection{5-Fold Cross-validation on NICE}
Table \ref{tab:nice-cv} reports 5-fold cross-validation results with mean and
standard deviation across folds.

\begin{table}[htbp]
\centering
\caption{Mean and standard deviation of 5-fold cross-validation results on the NICE multi-label dataset. Sentence-BERT achieves the highest Macro-F1 and LRAP; TF-IDF Linear SVM achieves the highest Micro-F1; the hybrid quantum--SVM model achieves the lowest Hamming loss.}
\label{tab:nice-cv}
\resizebox{\textwidth}{!}{
\begin{tabular}{lcccc}
\toprule
Model & Micro-F1 & Macro-F1 & Hamming Loss & LRAP \\
\midrule
Sentence-BERT + LR & 0.6493 $\pm$ 0.0357 & \textbf{0.6121 $\pm$ 0.0450} & 0.0919 $\pm$ 0.0200 & \textbf{0.8361 $\pm$ 0.0225} \\
TF-IDF Linear SVM & \textbf{0.6821 $\pm$ 0.0405} & 0.6049 $\pm$ 0.0336 & 0.0764 $\pm$ 0.0094 & 0.8056 $\pm$ 0.0262 \\
Hybrid Quantum--SVM Fusion & 0.6739 $\pm$ 0.0432 & 0.6000 $\pm$ 0.0318 & \textbf{0.0742 $\pm$ 0.0099} & 0.8075 $\pm$ 0.0259 \\
TF-IDF Logistic Regression & 0.6787 $\pm$ 0.0463 & 0.5977 $\pm$ 0.0229 & 0.0759 $\pm$ 0.0140 & 0.7990 $\pm$ 0.0240 \\
Contrastive Quantum Projection & 0.6145 $\pm$ 0.0493 & 0.5602 $\pm$ 0.0416 & 0.0873 $\pm$ 0.0116 & 0.7775 $\pm$ 0.0320 \\
Quantum-inspired Projection & 0.5931 $\pm$ 0.0166 & 0.5490 $\pm$ 0.0237 & 0.1096 $\pm$ 0.0213 & 0.7830 $\pm$ 0.0304 \\
Label-frequency baseline & 0.2328 $\pm$ 0.0034 & 0.2083 $\pm$ 0.0032 & 0.7927 $\pm$ 0.0019 & 0.4059 $\pm$ 0.0110 \\
\bottomrule
\end{tabular}
}
\end{table}

Under standard 5-fold cross-validation with a global threshold, Sentence-BERT
is the strongest Macro-F1 baseline. The hybrid quantum--SVM fusion model remains
competitive with TF-IDF baselines and improves over both pure quantum-inspired
variants. A Wilcoxon signed-rank test over fold-level Macro-F1 finds no
significant difference between the hybrid model and TF-IDF Linear SVM
($p=0.625$), TF-IDF Logistic Regression ($p=1.000$), or Sentence-BERT
($p=0.125$). Given the small number of folds and the dependence induced by
cross-validation resampling, these values are reported as descriptive evidence
rather than as a strong statistical conclusion.

\subsection{5-Fold Cross-validation with Per-label Thresholds}
Table \ref{tab:nice-per-label} reports the 5-fold results when each label
receives its own threshold selected on the validation split.

\begin{table}[htbp]
\centering
\caption{Mean and standard deviation of 5-fold results with per-label threshold calibration on the NICE multi-label dataset. Sentence-BERT has the highest Macro-F1 and LRAP, while the hybrid quantum--SVM model remains the strongest non-embedding model in Macro-F1.}
\label{tab:nice-per-label}
\resizebox{\textwidth}{!}{
\begin{tabular}{lcccc}
\toprule
Model & Micro-F1 & Macro-F1 & Hamming Loss & LRAP \\
\midrule
Sentence-BERT + LR & \textbf{0.6386 $\pm$ 0.0280} & \textbf{0.5990 $\pm$ 0.0276} & \textbf{0.0947 $\pm$ 0.0092} & \textbf{0.8361 $\pm$ 0.0225} \\
Hybrid Quantum--SVM Fusion & 0.6138 $\pm$ 0.0602 & 0.5775 $\pm$ 0.0520 & 0.1010 $\pm$ 0.0331 & 0.8075 $\pm$ 0.0259 \\
TF-IDF Linear SVM & 0.6085 $\pm$ 0.0588 & 0.5640 $\pm$ 0.0309 & 0.1057 $\pm$ 0.0378 & 0.8056 $\pm$ 0.0262 \\
TF-IDF Logistic Regression & 0.6219 $\pm$ 0.0588 & 0.5629 $\pm$ 0.0532 & 0.0967 $\pm$ 0.0195 & 0.7990 $\pm$ 0.0240 \\
\bottomrule
\end{tabular}
}
\end{table}

With per-label threshold calibration, Sentence-BERT again obtains the best
Macro-F1. The hybrid quantum--SVM model is the strongest non-embedding model in
Macro-F1, which suggests that label-specific thresholds help the hybrid model
handle imbalanced NFR labels.

\subsection{Ablation Study}
Table \ref{tab:nice-ablation} reports a 5-fold ablation study of the main
quantum-inspired components. The ablation compares the original positive
projection, contrastive projection with and without interference, an SVM-only
component using the same sublinear TF-IDF representation as the hybrid model,
and hybrid fusion with different quantum weights.

\begin{table}[htbp]
\centering
\caption{Ablation study on the NICE multi-label dataset. The best hybrid
configuration uses a small quantum weight ($\alpha=0.15$), indicating that the
projection score is most useful as a calibrated auxiliary signal rather than as
the dominant score.}
\label{tab:nice-ablation}
\resizebox{\textwidth}{!}{
\begin{tabular}{lcccc}
\toprule
Variant & Micro-F1 & Macro-F1 & Hamming Loss & LRAP \\
\midrule
Hybrid $\alpha=0.15$ & 0.6739 $\pm$ 0.0432 & 0.6000 $\pm$ 0.0318 & 0.0742 $\pm$ 0.0099 & 0.8075 $\pm$ 0.0259 \\
SVM-only sublinear TF-IDF & 0.6713 $\pm$ 0.0553 & 0.5945 $\pm$ 0.0409 & 0.0843 $\pm$ 0.0263 & 0.8081 $\pm$ 0.0291 \\
Hybrid $\alpha=0.50$ & 0.6501 $\pm$ 0.0544 & 0.5865 $\pm$ 0.0443 & 0.0818 $\pm$ 0.0119 & 0.7933 $\pm$ 0.0374 \\
Hybrid $\alpha=0.30$ & 0.6442 $\pm$ 0.0597 & 0.5761 $\pm$ 0.0470 & 0.0759 $\pm$ 0.0085 & 0.8029 $\pm$ 0.0294 \\
Contrastive projection with interference & 0.6145 $\pm$ 0.0493 & 0.5602 $\pm$ 0.0416 & 0.0873 $\pm$ 0.0116 & 0.7775 $\pm$ 0.0320 \\
Positive projection with interference & 0.5931 $\pm$ 0.0166 & 0.5490 $\pm$ 0.0237 & 0.1096 $\pm$ 0.0213 & 0.7830 $\pm$ 0.0304 \\
Contrastive projection without interference & 0.5995 $\pm$ 0.0353 & 0.5445 $\pm$ 0.0363 & 0.0878 $\pm$ 0.0089 & 0.7783 $\pm$ 0.0328 \\
\bottomrule
\end{tabular}
}
\end{table}

The ablation supports three observations. First, contrastive projection with
interference improves Macro-F1 over the original positive projection. Second,
the hybrid model improves over pure projection variants, showing the value of
discriminative calibration. Third, increasing the quantum weight beyond
$\alpha=0.15$ reduces Macro-F1, suggesting that the current projection signal
should complement rather than replace the SVM score.

\subsection{Explainability Evaluation}
The deletion test evaluates whether high-contribution explanation terms are
actually influential. Across 158 true label assignments in the held-out test
split, deleting the top-3 intrinsic projection terms reduced the mean projection
score from 0.6577 to 0.5413, a mean drop of 0.1165. Deleting three random
nonzero terms reduced the mean score to 0.6435, a mean drop of 0.0143. Thus,
targeted deletion produced a score drop about 8.17 times larger than random
deletion. For comparison, deleting the top-3 TF-IDF Linear SVM coefficient
features produced a mean score drop of 0.0867, about 4.09 times its random
deletion baseline. This supports RQ2 by showing that the highlighted projection
terms are not merely descriptive; they are also influential under the model's
scoring function and at least comparable to a standard linear-model feature
attribution baseline.

Table \ref{tab:qualitative} gives a qualitative example from the same deletion
test. The highest-contribution terms for the Security label correspond to access
control concepts that are also plausible human rationales.

\begin{table}[htbp]
\centering
\caption{Example requirement with top token-level contributions for the true NFR label, illustrating the intrinsic explanation produced by the projection model.}
\label{tab:qualitative}
\begin{tabular}{p{0.52\textwidth}p{0.16\textwidth}p{0.22\textwidth}}
\toprule
Requirement excerpt & Label & Top terms \\
\midrule
``The product shall ensure that it can only be accessed by authorized users'' & Security & authorized, only, access \\
\bottomrule
\end{tabular}
\end{table}

\subsection{Auxiliary PROMISE-expanded Result}
The auxiliary single-label PROMISE-expanded experiment shows that the
quantum-inspired projection model can be competitive in a single-label setting.
It obtains 0.7025 accuracy and 0.6748 Macro-F1, compared with 0.6899 accuracy
and 0.6344 Macro-F1 for TF-IDF Logistic Regression. However, this dataset is
not the main evidence for the multi-label research claim.

\section{Discussion}
The experiments provide a balanced picture. On the main NICE multi-label
dataset, Sentence-BERT is the strongest baseline in Macro-F1 and LRAP, which is
consistent with the expectation that semantic embeddings are strong baselines
for text classification. The hybrid quantum--SVM fusion model is competitive
with TF-IDF baselines and has the best Hamming loss under the global-threshold
protocol, but it does not dominate the embedding baseline.

The contrastive quantum variant performs better than the original projection
model in 5-fold cross-validation, showing that discriminative label directions
are more effective than positive-only centroids. The hybrid model further
improves stability by combining the interpretable projection signal with SVM
scores. The ablation study further shows that a small quantum fusion weight
works better than larger quantum weights in the current TF-IDF setting. The
strongest supported claim is therefore not that a pure
quantum-inspired model outperforms all classical baselines, but that a
calibrated quantum-inspired representation can be competitive while providing
token-level explanations that pass a deletion-based faithfulness check.

To reduce reliance on fold-level significance testing, we also ran a paired
bootstrap over the held-out NICE test requirements. The hybrid model's Macro-F1
advantage over TF-IDF Linear SVM was 0.0814 with a 95\% bootstrap interval of
[-0.0039, 0.1541], and its advantage over TF-IDF Logistic Regression was 0.0753
with interval [-0.0147, 0.1533]. These intervals overlap zero, so they support
the same conservative conclusion as the cross-validation results: the hybrid
model is competitive with TF-IDF baselines, but the evidence is not strong
enough to claim clear accuracy superiority.

Although Sentence-BERT obtains the highest Macro-F1 and LRAP, the
quantum-inspired hybrid model occupies a different point in the
accuracy--interpretability--cost trade-off. Sentence-BERT relies on a pretrained
transformer encoder and produces dense embeddings whose individual dimensions
are not directly interpretable, so explanations require post-hoc methods such
as LIME or SHAP \cite{ribeiro2016lime,lundberg2017shap}. In contrast, the
projection model yields token-level contributions that are intrinsic to the
scoring function and pass a deletion-based faithfulness check, while remaining
competitive with TF-IDF baselines and requiring no pretrained model. The method
is therefore most attractive when interpretability and a small dependency
footprint matter more than maximum Macro-F1.

\section{Threats to Validity}
\textbf{Dataset size.} NICE contains only a few hundred NFR-labeled
requirements after filtering. Some labels are rare, which makes evaluation
unstable and limits the power of fold-level statistical tests.

\textbf{Baseline scope.} The current implementation includes TF-IDF baselines
and Sentence-BERT, but not fine-tuned BERT, RoBERTa, NoRBERT, ML-kNN, or the
small-language-model classifier from the NICE paper.

\textbf{Model simplicity.} The quantum-inspired models are based on TF-IDF
semantic states, projection directions, and score fusion. They do not yet use
contextual neural embeddings inside the projection component.

\textbf{Threshold sensitivity.} Multi-label classification depends strongly on
decision thresholds. The experiments show that per-label threshold calibration
can change the relative ranking of models.

\textbf{Explainability evaluation.} The deletion test measures faithfulness to
the model's own scoring function, but it does not measure human usefulness or
agreement with expert rationales.

\textbf{Ablation scope.} The ablation study examines projection type,
interference, and fusion weight, but it does not yet test contextual embeddings
inside the quantum-inspired projection component.

\section{Conclusion}
This paper presented explainable quantum-inspired models for multi-label NFR
classification. Requirements are represented as semantic states, NFR labels are
modeled as projection directions, and predictions are explained through
semantic amplitude contributions. Experiments on the NICE multi-label dataset
show that the hybrid quantum--SVM model is competitive with TF-IDF baselines,
although Sentence-BERT is the strongest Macro-F1 and LRAP baseline. The
deletion-based explainability evaluation shows that top contribution terms are
substantially more influential than random terms. The ablation study indicates
that contrastive projection and light-weight hybrid fusion are the most useful
quantum-inspired components in the current implementation.

Future work should integrate contextual embeddings into the quantum-inspired
representation, compare with fine-tuned BERT/RoBERTa \cite{liu2019roberta} and
NoRBERT baselines, evaluate ML-kNN \cite{zhang2007mlknn} as a classical
multi-label baseline, and collect human rationale judgments for explanation
quality.

\section*{Reproducibility}
The implementation contains scripts for reproducing the experiments:
\begin{itemize}
    \item \texttt{scripts/run\_nice\_multilabel\_experiment.py}
    \item \texttt{scripts/run\_nice\_cv\_experiment.py}
    \item \texttt{scripts/run\_nice\_per\_label\_threshold\_experiment.py}
    \item \texttt{scripts/run\_nice\_ablation\_experiment.py}
    \item \texttt{scripts/run\_explainability\_deletion\_test.py}
    \item \texttt{scripts/run\_nice\_robustness\_experiment.py}
    \item \texttt{scripts/run\_all\_experiments.py}
    \item \texttt{scripts/run\_promise\_experiment.py}
\end{itemize}
For anonymous review, the artifact should be submitted as an anonymized
repository or archive together with the paper.

\bibliographystyle{splncs04}
\bibliography{references}

\end{document}
